% Experimental results of the two-tier model (English version of two_tier_results.tex, 2026-10-12).
% Requires: \usepackage{booktabs,multirow}
% Explanation and data sources: two_tier_results.md. The "TBD" cells of Table 2 are to be replaced when the
% validation readings are finished.

\begin{table}[htbp]
  \centering
  \caption{Single-aircraft closed-loop generation on the validation split of the labelled data set}
  \label{tab:two_tier_val_single}
  \begin{tabular}{lrrcc}
    \toprule
    \multirow{2}{*}{Model} & \multirow{2}{*}{Trajectories} & \multirow{2}{*}{Landing rate} & arc ADE (m) & arc FDE (m) \\
     & & & mean / median & mean / median \\
    \midrule
    Labelled instructions, closed-loop replay (reference) & 947   & 100.0\% & 83 / 69      & 132 / 103 \\
    \midrule
    Base model                                            & 1\,894 & 94.6\%  & 1\,444 / 315 & 775 / 192 \\
    Stage C, round 5 (C-r5)                               & 1\,894 & 97.9\%  & 1\,242 / 298 & 708 / 190 \\
    Stage D, round 7 (D-r7)                               & 1\,894 & 96.9\%  & 1\,248 / 295 & 674 / 195 \\
    \bottomrule
  \end{tabular}
  \par\smallskip
  \parbox{0.95\linewidth}{\footnotesize
    200 validation flights are drawn at random at each of the 5 airports (the same flights for all three models), and
    each flight is generated twice. Only flights in the prior's selection are counted: those whose labelled outcome is a
    landing and whose observed track has no fault, 947 flights in all. arc ADE: the generated and the observed
    trajectory are each parametrised by normalised horizontal arc length, and the 3D distance is averaged at the arc
    fractions $k/64$ ($k=1,\dots,64$). arc FDE: the 3D distance at arc fraction 1 (the end point). The reference row is
    the error of the labelled instruction sentences, replayed through the executor, against the observed track.}
\end{table}

\begin{table}[htbp]
  \centering
  \caption{Multi-aircraft closed-loop control on the traffic windows of the validation split (post-training)}
  \label{tab:two_tier_val_multi}
  \begin{tabular}{lrcc}
    \toprule
    Model & Controlled aircraft & Landing rate & Loss-of-separation rate \\
    \midrule
    Base model              & TBD & TBD & TBD \\
    Stage C, round 5 (C-r5) & TBD & TBD & TBD \\
    Stage D, round 7 (D-r7) & TBD & TBD & TBD \\
    \bottomrule
  \end{tabular}
  \par\smallskip
  \parbox{0.95\linewidth}{\footnotesize
    Multi-aircraft traffic windows of stage D (durations 5, 10 and 20\,min; 100 real windows for each airport and
    duration); the three models use the same windows and the same random numbers. The denominator of each rate is the
    number of controlled aircraft. Loss of separation: the distance between a controlled aircraft and another aircraft
    falls below the separation rule and the controlled aircraft is responsible for it. The base model runs with the
    output of the traffic module set to zero.}
\end{table}
